\hspace{3mm}

\centerline{\bf \Huge Инварианты}

\centerline{\bf \large Базовый уровень}

\begin{flushright}
        \scriptsize{}
\end{flushright}

\problem{0.1} 

\problem{0.2} 


\problem{1}
Можно ли доску 8 × 8 разрезать по границам клеток на части и сложить из этих частей клетчатый прямоугольник 5 × 13?

\problem{2}
На доске написано 10 плюсов и 15 минусов. Разрешается стереть любые два знака и написать вместо них плюс, если они одинаковы, и минус в противном случае. Какой знак останется на доске после выполнения 24 таких операций?


\problem{3}
На доске написаны числа $1, 2, 3, ..., 19, 20$. Разрешается стереть любые два числа $a$ и $b$ и вместо них написать число $a + b$. Какое число может остаться на доске после 19 таких операций?

\problem{4}
Разменный автомат меняет одну монету на пять других. Можно ли с его помощью разменять одну монету на 100 монет?

\problem{5}
Дана доска $8 \times 8$, раскрашенная в шахматном порядке. За одно действие можно выбрать любые две соседние клетки и перекрасить их в противоположные цвета: белые в черный, а черные в белый. Можно ли за несколько действий оставить на доске ровно одну черную клетку?

\problem{6}
На доске записано несколько нулей, единиц и двоек. Разрешается стереть две неравные цифры и записать вместо них одну цифру, отличную от стёртых. Докажите, что если в результате нескольких таких операций на доске останется одна-единственная цифра, то она не зависит от порядка, в котором производились стирания.



\newpage

\hspace{3mm}

\centerline{\bf \Huge Инварианты}

\centerline{\bf \large Продвинутый уровень}

\begin{flushright}
        \scriptsize{}
\end{flushright}

\problem{1}
$100$ фишек выставлены в ряд. Разрешено менять местами две фишки, стоящие через одну фишку. Можно ли с помощью таких операций переставить все фишки в обратном порядке?

\problem{2}
На доске написаны числа $1, 2, 3, ..., 19, 20$. Разрешается стереть любые два числа $a$ и $b$ и вместо них написать число $a + b - 1$. Какое число может остаться на доске после 19 таких операций?

\problem{3} 
У каждого депутата в парламенте не более трёх врагов (вражда обоюдна). Пар- ламент разбит на 2 палаты. Каждый день один из депутатов, обнаружив в своей палате не менее двух врагов, переходит в другую палату. Докажите, что рано или поздно переходы прекратятся.

\problem{4}
Даны три кучки камней, по $n$ камней в каждой. За один ход можно выбрать две кучки, убрать из них по одному камню, при этом добавив один камень в третью кучку. При каких $n$ можно через несколько ходов оставить только один камень?

\problem{5} 
На доске написаны три числа. Каждую минуту они изменяются по следую- щему правилу: если на данный момент на доске находятся числа $x, y, z$ то они меняются на $2x - 2y + z -1, 2y - 2z + x - 1, 2z - 2x + y - 1$ соответственно. Докажите, что когда-нибудь на доске появится отрицательное число.

\problem{6}
На доске выписаны числа $1, \dfrac{1}{2}, \dfrac{1}{3}, ..., \dfrac{1}{100}.$ Выбираем из написанных на доске два произвольных числа $a$ и $b$, стираем их и пишем на доску число $a + b + ab.$ Такую операцию проделываем 99 раз, пока не останется одно число. Какое это число? Найдите его и докажите, что оно не зависит от последовательности выбора чисел.

